\chapter{The Rewrite Engine}\label{ch:rewrite}

A dataplane is assembled from atoms and then kept in canonical shape by the
rewrite engine: the runtime's equational theories, directed so that the
engine applies them in one direction only, replacing redundant or expensive
subexpressions by cheaper ones and reducing every expression to a unique
normal form. This chapter formalizes that machinery. The equational theories
of \autoref{ch:lawvere}: the stack theory of LIFO buffers and its pointer
fragment (\autoref{thm:9}--\autoref{thm:12}), the parser theory of sequence and choice
(\autoref{thm:10}), and the transport theory of batched I/O
(\autoref{thm:11}): become term rewrite systems; monoidal coherence is
absorbed into the representation by Mac Lane's strictification, so coherence
generates no rewrite rules of its own \cite{maclane}. Termination is proved
by well-founded orders (lexicographic path orders), confluence follows from
critical-pair analysis and Newman's lemma \cite{newman}, and the systems are
produced by Knuth--Bendix completion, the term-rewriting analogue of
computing a Gr\"obner basis \cite{knuthbendix}. The seven theorems
\autoref{thm:13}--\autoref{thm:19}, stated and proved in full below, then establish: termination,
local confluence, and completion for the ring system (\autoref{thm:13}--\autoref{thm:15});
confluence and unique normal forms for the parser and batch rewrites
(\autoref{thm:16}--\autoref{thm:17}); soundness of normalization for behavior and cost (\autoref{thm:18}); and
decidability of the full runtime equational theory on finite signatures
(\autoref{thm:19}).

\paragraph{Term rewrite systems.}
Fix a \emph{signature}: a finite set $\Sigma$ of operation symbols, each with
a fixed arity, and a countably infinite set $X$ of variables. \emph{Terms}
are built as usual: variables are terms, and for an $n$-ary symbol $f$ and
terms $t_1, \ldots, t_n$ the expression $f (t_1, \ldots, t_n)$ is a term; the
\emph{size} $|t|$ counts its operation symbols. A \emph{substitution}
$\sigma$ maps variables to terms, extended homomorphically; a \emph{rewrite
rule} is a pair $l \to r$ of terms; and a \emph{term rewrite system} (TRS)
is a set of rules. A \emph{rewrite step} $C[l\sigma] \to_R C[r\sigma]$
replaces an instance of a left-hand side inside any context $C$, and
$\to_R^{*}$ is its reflexive--transitive closure. A term is a \emph{normal
form} when no step applies to it, and $\mathrm{NF} (t)$ is its normal form
when one exists and is unique.

\paragraph{Termination, confluence, critical pairs.}
A TRS is \emph{terminating} if it admits no infinite rewrite sequence. A
\emph{reduction order} is a well-founded order $\succ$ on terms, closed
under contexts and substitutions, with $l \succ r$ for every rule $l \to r$;
any reduction order proves termination. The order used below is the
classical \emph{lexicographic path order} (LPO): fix a \emph{precedence} on
the operation symbols and compare $f (s_1, \ldots, s_m)$ with
$g (t_1, \ldots, t_n)$ by (i) $s_i \succeq t$ for some $i$, or (ii) $f$
higher in precedence than $g$ and the left term exceeds every $t_j$, or
(iii) $f = g$ and the argument tuples compare lexicographically; LPO is a
reduction order over finite signatures. A TRS is \emph{confluent} if any
two terms reachable from a common term are \emph{joinable} (both reduce to
a common term), and \emph{locally confluent} if the same holds for one-step
reachability. A \emph{critical pair} is a pair of one-step reducts produced
by overlapping two left-hand sides at a non-variable position and unifying
at their most general overlap; the classical Knuth--Bendix criterion says a
TRS is locally confluent iff every critical pair is joinable
\cite{knuthbendix}. Termination plus local confluence yields confluence by
Newman's lemma, proved below \cite{newman}.

\paragraph{The three systems.}
The engine normalizes three theories on disjoint sorts, so rules of
different theories never overlap. The \emph{ring system} $R$ has the two
cancellation rules
\[
\mathit{push}\bigl (\mathit{pop} (x)\bigr) \to x, \qquad
\mathit{pop}\bigl (\mathit{push} (x)\bigr) \to x ,
\]
which read, in the pipeline notation of \autoref{thm:9}, as
$\mathit{push};\mathit{pop} = \mathrm{id}$ and
$\mathit{pop};\mathit{push} = \mathrm{id}$ on stacks (and on the pointer
line $\mathbb{Z}$). The non-emptiness guard on
the second equation of \autoref{thm:9} is satisfied by the pattern itself:
$\mathit{pop} (\mathit{push} (x))$ denotes a buffer that has just received
the element pushed by the inner $\mathit{push}$ and is non-empty by
construction. The \emph{parser system} $P$ has the single left-factoring
rule
\[
(x \cdot a) \oplus (x \cdot b) \to x \cdot (a \oplus b) ,
\]
with $\oplus$ ordered choice and $\cdot$ sequence; the runtime stores every
alternative as a sequence, so an atom is a sequence of length one and the
rule factors atom prefixes as well. The \emph{batch system} $B$ has the
single flattening rule
\[
\mathrm{batch}\bigl (\mathrm{batch} (x,y), z\bigr) \to
\mathrm{batch}\bigl (x, \mathrm{batch} (y,z)\bigr),
\]
the binary restriction of the $n$-ary batch of \autoref{thm:46}.

\paragraph{Newman's lemma.}
A classical result bridges local to global confluence: a terminating and
locally confluent rewrite relation is confluent \cite{newman}. Since $\to$
is terminating, its strict transitive closure $\to^{+}$ is well-founded, and
we argue by well-founded induction on $t$ that any two terms reachable from
$t$ are joinable. If $t$ is a normal form, the only term reachable from $t$
is $t$ itself. Otherwise let $t \to^{*} u$ and $t \to^{*} v$; if either
sequence is empty, $u$ and $v$ are trivially joinable, so assume both begin
with one step, $t \to t_1 \to^{*} u$ and $t \to t_2 \to^{*} v$. By local
confluence $t_1$ and $t_2$ are joinable: $t_1 \to^{*} w$ and
$t_2 \to^{*} w$. The terms $t_1$, $t_2$, and $w$ are proper reducts of $t$,
so the induction hypothesis applies to each: $u$ and $w$ are joinable and
$v$ and $w$ are joinable, yielding $u \to^{*} w_1$, $w \to^{*} w_1$,
$v \to^{*} w_2$, $w \to^{*} w_2$; and since $w_1$ and $w_2$ are reachable
from $w$, they are joinable, $w_1 \to^{*} z$, $w_2 \to^{*} z$. Then
$u \to^{*} z$ and $v \to^{*} z$: $u$ and $v$ join, and induction gives
confluence.

\setcatalog{13}
\begin{theorem}[The ring rewrite system is terminating (\autoref{thm:13})]\label{thm:13}
\textbf{[N].} The ring rewrite system $\mathit{push} (\mathit{pop} (x)) \to x$
and $\mathit{pop} (\mathit{push} (x)) \to x$ is terminating: in the
lexicographic path order with precedence
$\mathit{push} \succ \mathit{pop}$ every rule strictly decreases the term,
and every rule also strictly decreases the number of operation symbols.
\end{theorem}

\begin{proof}
We verify the lexicographic path order directly. For the first rule,
$\mathit{push} (\mathit{pop} (x)) \succ x$: the immediate subterm
$\mathit{pop} (x)$ satisfies $\mathit{pop} (x) \succ x$, since
$\mathit{pop} (x)$ has the immediate subterm $x$ with $x \succeq x$ (clause
(i)), and clause (i) applied to $\mathit{push} (\mathit{pop} (x))$ gives the
claim. Symmetrically $\mathit{pop} (\mathit{push} (x)) \succ x$, its
immediate subterm $\mathit{push} (x)$ dominating $x$ the same way. LPO is a
reduction order over the finite signature $\{\mathit{push}, \mathit{pop}\}$,
so every step decreases a well-founded order and the system is terminating.
The size witness is immediate: each left-hand side has two more operation
symbols than its right-hand side, so $|\mathit{push} (\mathit{pop} (x))| =
|x| + 2 > |x|$, and similarly for the second rule.
\end{proof}

\setcatalog{14}
\begin{theorem}[Local confluence of the ring system (\autoref{thm:14})]\label{thm:14}
\textbf{[N].} The ring rewrite system of \autoref{thm:13} is locally
confluent: every critical pair is joinable. It is therefore confluent, by
Newman's lemma, since it is terminating \cite{newman}.
\end{theorem}

\begin{proof}
Let $l_1 = \mathit{push} (\mathit{pop} (x))$ and
$l_2 = \mathit{pop} (\mathit{push} (y))$ (variables renamed apart). The
non-variable positions of $l_1$ are the root and the argument of
$\mathit{push}$, holding the subterm $\mathit{pop} (x)$. At the argument
position, $\mathit{pop} (x)$ unifies with $l_2$ under
$\theta = \{x \mapsto \mathit{push} (y)\}$, and the two one-step rewrites of
the overlap term $\mathit{push} (\mathit{pop} (\mathit{push} (y)))$ both
produce $\mathit{push} (y)$: the root redex substitutes $\theta$ into the
right-hand side, and the inner redex $\mathit{pop} (\mathit{push} (y))$,
reduced by $l_2$, leaves $\mathit{push} (y)$ as well. Symmetrically, the
subterm $\mathit{push} (y)$ at the argument position of $l_2$ unifies with
$l_1$ under $\{y \mapsto \mathit{pop} (x)\}$, and the two one-step rewrites
of $\mathit{pop} (\mathit{push} (\mathit{pop} (x)))$ both produce
$\mathit{pop} (x)$. At the root the left-hand sides do not unify, since
$\mathit{push} \neq \mathit{pop}$, and the remaining positions hold
variables, where redexes never occur. Hence the only critical pairs are the
trivial root self-overlaps and the pairs
$\langle \mathit{push} (y), \mathit{push} (y) \rangle$ and
$\langle \mathit{pop} (x), \mathit{pop} (x) \rangle$, each joining in a
single step, so every critical pair is joinable and the Knuth--Bendix
criterion gives local confluence \cite{knuthbendix}; with termination
(\autoref{thm:13}), Newman's lemma gives confluence \cite{newman}. The
normal forms are the terms $\mathit{push}^{k} (x)$ and $\mathit{pop}^{k} (x)$,
$k \geq 0$: a term that avoids both patterns
$\mathit{push} (\mathit{pop} (\cdot))$ and
$\mathit{pop} (\mathit{push} (\cdot))$ has all its operation symbols equal.
\end{proof}

\setcatalog{15}
\begin{theorem}[Knuth--Bendix completion of the ring theory (\autoref{thm:15})]\label{thm:15}
\textbf{[N].} Knuth--Bendix completion of the ring theory, with the
equations of \autoref{thm:9} oriented by the size-lexicographic order,
terminates in a complete system: it returns exactly the two cancellation
rules of \autoref{thm:13}, which are terminating and confluent. The
computational claim is script-verified: the checker
\texttt{kb\_completion} of \autoref{app:a} executes the completion, and its
run prints \texttt{PASS} \cite{knuthbendix}.
\end{theorem}

\begin{proof}
The completion procedure orients every equation by a well-founded order:
here the size-lexicographic order (smaller size first, ties broken
lexicographically): then computes critical pairs and adds a rule for
every pair that fails to join, until all critical pairs join. The equations
$\mathit{push} (\mathit{pop} (x)) = x$ and
$\mathit{pop} (\mathit{push} (x)) = x$ are already oriented: each left-hand
side has size $2$ and the right-hand side size $0$, so the initial rules are
exactly $R_1$ and $R_2$ of \autoref{thm:13}. By the critical-pair computation
of \autoref{thm:14}, the two nontrivial critical pairs join in one step and
the trivial self-overlaps join in zero steps; no critical pair fails to
join, so no rule is ever added and the procedure terminates after the first
pass. The resulting system is terminating (\autoref{thm:13}) and confluent
(\autoref{thm:14}), hence complete, and because the rules are exactly the
equations, it presents the same theory and decides it: two terms are equal
in the ring theory iff they reduce to a common term, iff their normal forms
coincide. The script \texttt{kb\_completion} implements terms, matching,
rewriting, normalization, critical-pair computation, and the completion loop
for these two equations; it asserts that completion terminates, that every
critical pair joins, that the original rules are present, and that the
system decides the theory on sample terms, and it prints \texttt{PASS}. The
run is recorded in \autoref{app:a} \cite{knuthbendix}.
\end{proof}

\setcatalog{16}
\begin{theorem}[Parser left factoring (\autoref{thm:16})]\label{thm:16}
\textbf{[N].} The parser rewrite system with the single left-factoring rule
$(x \cdot a) \oplus (x \cdot b) \to x \cdot (a \oplus b)$ is confluent and
terminating. Consequently every parser term has a unique left-factored
normal form: a term in which no choice node has two branches sharing a
common first component.
\end{theorem}

\begin{proof}
Termination: for any substitution $\sigma$, the left-hand side
$(x \cdot a) \oplus (x \cdot b)$ contains $2|x\sigma| + |a\sigma| +
|b\sigma| + 3$ operation symbols and the right-hand side
$x \cdot (a \oplus b)$ contains $|x\sigma| + |a\sigma| + |b\sigma| + 2$;
the difference $|x\sigma| + 1 \geq 1$ is positive, so every step decreases
the size and no infinite sequence exists. Local confluence: the left-hand
side $l = (x \cdot a) \oplus (x \cdot b)$ has non-variable positions
$\varepsilon$, $1$, and $2$; the subterms at positions $1$ and $2$ are
$\cdot$-terms, and no $\cdot$-term unifies with $l$, whose root is $\oplus$,
so the only overlap of $l$ with a renaming of itself is the trivial root
overlap, joining in zero steps. With no other rules, every critical pair is
joinable, the system is locally confluent, and Newman's lemma gives
confluence; termination then yields a unique normal form for every term: a
term with no redex, i.e.\ no choice node $\oplus$ whose two branches are
sequences with identical first component, which is the left-factored form.
The rule is sound by \autoref{thm:10} (sequence distributes over choice),
and since the runtime stores each alternative as a sequence, atom branches
are sequences of length one and atom prefixes are factored as well.
\end{proof}

\setcatalog{17}
\begin{theorem}[Batch-flattening rewrites (\autoref{thm:17})]\label{thm:17}
\textbf{[N].} The batch rewrite system with the single rule
$\mathrm{batch} (\mathrm{batch} (x,y), z) \to
\mathrm{batch} (x, \mathrm{batch} (y,z))$ is terminating and confluent on
batch trees. Consequently every batch tree has a unique flattened normal
form: a right spine $\mathrm{batch} (x_1, \mathrm{batch} (x_2, \ldots,
\mathrm{batch} (x_{n-1}, x_n) \ldots))$, identified with the single $n$-ary
batch of \autoref{thm:46}.
\end{theorem}

\begin{proof}
Termination: in the lexicographic path order with the single operation
symbol $\mathrm{batch}$ and lexicographic argument status, the rule
decreases, since the first arguments compare
$\mathrm{batch} (x,y) \succ x$ (clause (i): the immediate subterm $x$
satisfies $x \succeq x$), and the lexicographic comparison is decided by the
first argument. LPO is a reduction order over the finite signature
$\{\mathrm{batch}\}$, so the system is terminating. Local confluence: the
left-hand side $l = \mathrm{batch} (\mathrm{batch} (x,y), z)$ has
non-variable positions $\varepsilon$ and $1$; the subterm at position $1$,
$\mathrm{batch} (x,y)$, unifies with a renaming of $l$ under
$x \mapsto \mathrm{batch} (x',y')$, $y \mapsto z'$, producing the overlap
term $\mathrm{batch} (\mathrm{batch} (\mathrm{batch} (x',y'),z'), z)$. The
outer redex rewrites this to $\mathrm{batch} (\mathrm{batch} (x',y'),
\mathrm{batch} (z',z))$, and one further step brings it to
$\mathrm{batch} (x', \mathrm{batch} (y', \mathrm{batch} (z',z)))$; the inner
redex rewrites the overlap term to
$\mathrm{batch} (\mathrm{batch} (x', \mathrm{batch} (y',z')), z)$, and two
further steps reach the same term. The critical pair joins, and the root
self-overlap is trivial; the Knuth--Bendix criterion gives local
confluence, and Newman's lemma gives confluence. The unique normal form is
the term with no redex, i.e.\ a right spine of $\mathrm{batch}$
applications, the flattened batch of \autoref{thm:46}.
\end{proof}

\setcatalog{18}
\begin{theorem}[Normalization soundness (\autoref{thm:18})]\label{thm:18}
\textbf{[N].} Suppose every rule of the three systems is sound: the
equation it encodes holds as behavioral equivalence of molecules: and
cost-nonincreasing: $M \to N$ implies $\mathrm{cost} (N) \leq
\mathrm{cost} (M)$, componentwise on the cost vector of
\autoref{ch:enrich}. Then normalization is sound and cost-monotone: for
every molecule $M$,
\[
\mathrm{NF} (M) \approx M
\qquad\text{and}\qquad
\mathrm{cost}\bigl (\mathrm{NF} (M)\bigr) \leq \mathrm{cost} (M) .
\]
\end{theorem}

\begin{proof}
By termination and confluence of the three systems (\autoref{thm:13}--\autoref{thm:17}), every term
has a unique normal form and $\mathrm{NF}$ is a function: it is the common
endpoint of all maximal reduction sequences from the input. For soundness,
each rule is behavior-preserving by hypothesis: concretely, the ring
rules are the equations of \autoref{thm:9}, the parser rule is the
distributive law of \autoref{thm:10}, and the batch rule is the flattening
identity of \autoref{thm:46}: so a single step satisfies $M \approx N$;
behavioral equivalence is an equivalence relation (\autoref{ch:mol}), and an
induction on the length of the reduction sequence from $M$ to
$\mathrm{NF} (M)$ gives $\mathrm{NF} (M) \approx M$. For cost, a single step
satisfies $\mathrm{cost} (N) \leq \mathrm{cost} (M)$ in every component by
hypothesis, and chaining the inequalities along the sequence, by induction
on its length, gives $\mathrm{cost} (\mathrm{NF} (M)) \leq \mathrm{cost} (M)$
componentwise. The two halves are independent; the concrete rules satisfy
both hypotheses: cancellation removes a redundant pair of operations,
factoring parses a shared prefix once instead of once per branch, and
flattening replaces a nested batch by a flat one with the same I/O count.
\end{proof}

\setcatalog{19}
\begin{theorem}[Decidability of the runtime equational theory (\autoref{thm:19})]\label{thm:19}
\textbf{[N].} Let $E$ be the full runtime equational theory: the union of
the ring theory (\autoref{thm:13}--\autoref{thm:15}), the parser theory
(\autoref{thm:16}), and the transport theory, whose rewrite rules are the
batch-flattening rewrites of \autoref{thm:17}, presented over a finite
signature in which the three theories occupy disjoint sorts. Then $E$ is
decidable via normalization: for terms $s$ and $t$ of the same sort,
\[
s =_E t \iff \mathrm{NF} (s) = \mathrm{NF} (t),
\]
and both normal forms are effectively computable.
\end{theorem}

\begin{proof}
Within each sort the theory is presented by the rules that
\autoref{thm:13}--\autoref{thm:17} prove terminating and confluent, so the
rewrite relation and the equational theory coincide: $s =_E t$ iff $s$ and
$t$ rewrite, in both directions, to a common term, iff $\mathrm{NF} (s) =
\mathrm{NF} (t)$, the last equivalence by confluence. The systems act on
disjoint sorts, so no rule of one contains a symbol of another and redexes
never overlap. Termination of the combined engine: each component
normalizes every term of its sort in finitely many steps, and steps of one
theory leave the other sorts' terms untouched; any infinite run would
therefore contain infinitely many steps of a single sort, each strictly
decreasing that sort's well-founded size measure while the others stay
constant, a contradiction. Confluence: every maximal reduction sequence
normalizes each sort component, and each component has a unique normal
form, so all maximal sequences from the same term end at the same mixed
normal form: $\mathrm{NF}$ is a function. Decidability: the signature is
finite, so terms are finite trees over a finite alphabet, every step is an
effective replacement, and normalization terminates, so $\mathrm{NF} (s)$
and $\mathrm{NF} (t)$ are computable in finite time, and comparing them
decides $s =_E t$. Coherence contributes no equations: by Mac Lane's
coherence theorem and strictification the tensor is strictly associative
with a strict unit, and the braiding is a morphism rather than a rewrite,
so the coherence laws hold in every model without entering $E$
\cite{maclane}.
\end{proof}

\paragraph{Typed composition trees and multicategories.}
The expressions the engine rewrites are embedded in molecules: a molecule
over an atom signature $\Sigma$ is a typed composition tree whose leaves
are atoms and whose internal nodes are applications of composition $\circ$
and tensor $\otimes$, each node carrying the types of its domain and
codomain. Such trees, with the grafting substitution that plugs the output
of one tree into an input position of another, form a \emph{multicategory}:
objects are the types, and a multimorphism is a tree with a list of typed
inputs and one typed output \cite{leinster}. The rewrite rules are exactly
the identities this substitution structure respects: batch flattening is
the associativity of substitution, parser factoring the distributivity of
sequence over choice, ring cancellation the mutual inverseness of the ring
operations: and Kelly's \emph{clubs} axiomatize the general situation: a
club is a category of typed trees whose substitution structure presents the
free algebras of a 2-monad, and the free symmetric monoidal category
$F (\Sigma)$ of \autoref{ch:free} is presented by the club of permutations
acting on the atoms \cite{kelly, leinster}. The engine is therefore not an
ad-hoc optimizer but the computational reflection of the multicategorical
structure of Mol, and the normal forms of \autoref{thm:13}--\autoref{thm:19} are the canonical
forms of typed composition trees, whose equality \autoref{thm:19} makes decidable.
\section{Extending the Engine: Modular Completion}\label{sec:rewrite-modular}

A dataplane engine is never closed: new protocols and new atoms adjoin
new equational theories (new rewrite rules) to the systems of
\autoref{ch:lawvere}. The value of the completed systems
\ref{thm:13}--\ref{thm:19} is that they are \emph{stable under extension}:
this section proves the conditions under which adjoining a theory
preserves termination and confluence, so that the normal forms, the
decision procedure \ref{thm:19}, and the lookup table of \autoref{ch:free}
survive the addition. Confluence of a disjoint union is Toyama's
modularity theorem \citep{toyama}; termination of a disjoint union is
not modular in general, and is recovered here from simple/LPO
termination of the components.

\begin{definition}[Disjoint signatures and commuting systems]\label{def:modular}
Two rewrite systems $(\Sigma_1, E_1, \to_1)$ and $(\Sigma_2, E_2, \to_2)$
have \emph{disjoint signatures} when $\Sigma_1$ and $\Sigma_2$ share no
function symbols. They \emph{commute} when, for all terms $t, u, v$ with
$t \to_1 u$ and $t \to_2 v$, there exists $w$ with $u \to_2^* w$ and
$v \to_1^* w$. A sufficient condition for commutation is that every
$\to_2$-rule has a $\to_1$-normal left- and right-hand side and every
$\to_1$-rule has a $\to_2$-normal left- and right-hand side, in a
non-deleting signature: no rule discards proper subterms, so redexes of
the two systems never interfere. The \emph{componentwise normal form}
of a term under a disjoint union is the term obtained by reducing each
$\Sigma_i$-colored subterm to its $\to_i$-normal form.
\end{definition}

\setcatalog{54}
\begin{theorem}[Modular completion, \autoref{thm:54}]\label{thm:54}
Let $(\Sigma_1, E_1, \to_1)$ and $(\Sigma_2, E_2, \to_2)$ be rewrite
systems, each locally confluent.
\begin{enumerate}
\item[(i)] \emph{Confluence is modular for disjoint signatures}
\citep{toyama}. If $\Sigma_1 \cap \Sigma_2 = \emptyset$ and each system
is confluent, then $\to_1 \cup \to_2$ is confluent.
\item[(ii)] \emph{Termination is not modular in general.} There exist
disjoint terminating systems whose union does not terminate
\citep{toyama}. If, additionally, each system is simply terminating
(admitted by a simplification order, in particular by an LPO), then the
union is simply terminating under any combined precedence that extends
both component precedences, and is therefore terminating.
\item[(iii)] Consequently, if the signatures are disjoint, each system
is LPO-terminating and locally confluent, the union is terminating and
confluent, hence complete for $E_1 \cup E_2$; the normal form of a term
is its componentwise normal form, and the union decides $E_1 \cup E_2$.
The three runtime systems of this chapter (stack, parser, batch) satisfy
the hypotheses, so their union is complete.
\end{enumerate}
\end{theorem}

\begin{proof}
\emph{(i).} This is Toyama's modularity theorem for confluence: the
disjoint union of confluent term rewrite systems is confluent
\citep{toyama}. We record the statement rather than reproducing the
full argument; the engine uses it as a black box for the three
runtime systems, whose signatures $\{\mathit{push},\mathit{pop}\}$,
$\{\cdot,\oplus\}$, and $\{\mathrm{batch}\}$ are pairwise disjoint.

\emph{(ii).} Toyama also showed that termination is not a modular
property: there are disjoint terminating systems whose union admits an
infinite reduction \citep{toyama}. Simple termination is modular. A
reduction order that is a simplification order (closed under contexts
and substitutions, and containing the embedding relation) remains a
simplification order on the disjoint union of signatures when the
precedence is the disjoint union of the component precedences. LPO is
a simplification order. If each system is LPO-terminating, the union
is LPO-terminating under the combined precedence, hence terminating.

\emph{(iii).} The stack, parser, and batch systems are pairwise
signature-disjoint. Each is LPO-terminating (\autoref{thm:13},
\autoref{thm:16}, \autoref{thm:17}) and locally confluent, hence
confluent by Newman's lemma. By (ii) the union is terminating; by (i)
it is confluent. A terminating confluent system is complete for its
equational theory: $s =_{E_1 \cup E_2} t$ if and only if
$\mathrm{NF} (s) = \mathrm{NF} (t)$. By disjointness a normal form of
the union is a term with no $\to_1$-redex and no $\to_2$-redex, the
componentwise normal form. Each component normal form is computable
because each component theory is decidable (\autoref{thm:19}), so the
union decides $E_1 \cup E_2$.
\end{proof}

\begin{remark}
The engine's extension contract is therefore LPO-termination plus
signature disjointness (or, more generally, simple termination plus
Toyama confluence), not bare termination of the components. Adding a
theory whose rules are LPO-decreasing on a disjoint signature
preserves the normal forms, the decision procedure, and the lookup
table, so the guarantees of \ref{thm:13}--\ref{thm:19} and
\autoref{ch:free} hold for the extended system.
\end{remark}
